%======================================================================
% preamble.tex
%
% Shared preamble for
%   "Mathematics for Machine Learning and Deep Learning:
%    From Absolute Beginner to Prerequisites"
%
% Loaded by main.tex via %======================================================================
% preamble.tex
%
% Shared preamble for
%   "Mathematics for Machine Learning and Deep Learning:
%    From Absolute Beginner to Prerequisites"
%
% Loaded by main.tex via %======================================================================
% preamble.tex
%
% Shared preamble for
%   "Mathematics for Machine Learning and Deep Learning:
%    From Absolute Beginner to Prerequisites"
%
% Loaded by main.tex via %======================================================================
% preamble.tex
%
% Shared preamble for
%   "Mathematics for Machine Learning and Deep Learning:
%    From Absolute Beginner to Prerequisites"
%
% Loaded by main.tex via \input{preamble}. Everything that controls the
% look of the notes (fonts, colours, boxes, headings), the notation
% macros, and the automatic exercise-solution machinery lives here, so
% that every chapter file contains only content.
%======================================================================

%----------------------------------------------------------------------
% Fonts, encoding and micro-typography
%----------------------------------------------------------------------
\usepackage[T1]{fontenc}
\usepackage{lmodern}
\usepackage[final]{microtype}

%----------------------------------------------------------------------
% Page geometry and paragraph layout
%----------------------------------------------------------------------
\usepackage[a4paper,top=27mm,bottom=27mm,left=30mm,right=30mm,
            headheight=14pt]{geometry}
\usepackage[skip=5pt plus 1pt minus 1pt,indent=0pt]{parskip}
\usepackage{needspace}

%----------------------------------------------------------------------
% Mathematics
%----------------------------------------------------------------------
\usepackage{amsmath}
\usepackage{amssymb}
\usepackage{amsthm}
\usepackage{mathtools}
\usepackage{bm}

%----------------------------------------------------------------------
% Colours (one colour per kind of box, used consistently throughout)
%----------------------------------------------------------------------
\usepackage{xcolor}
\definecolor{MLchapter}{RGB}{25,60,110}      % chapter and section titles
\definecolor{MLlink}{RGB}{20,80,160}         % internal hyperlinks
\definecolor{MLurl}{RGB}{140,40,110}         % external hyperlinks
\definecolor{MLterm}{RGB}{20,55,120}         % newly introduced terms
\definecolor{MLdef}{RGB}{31,95,168}          % definitions
\definecolor{MLthm}{RGB}{128,40,110}         % theorems, propositions
\definecolor{MLexample}{RGB}{40,120,60}      % worked examples
\definecolor{MLintuition}{RGB}{176,104,0}    % intuition / picture this
\definecolor{MLml}{RGB}{0,118,124}           % why this matters for ML
\definecolor{MLpitfall}{RGB}{185,45,40}      % common pitfalls
\definecolor{MLremark}{RGB}{100,100,100}     % remarks
\definecolor{MLmotivation}{RGB}{70,60,160}   % motivation
\definecolor{MLnotation}{RGB}{70,100,115}    % notation
\definecolor{MLexercise}{RGB}{60,60,60}      % exercises
\definecolor{MLalgorithm}{RGB}{120,80,40}    % algorithms
\definecolor{MLsummary}{RGB}{25,60,110}      % chapter summaries
% Colours used inside figures
\definecolor{PlotBlue}{RGB}{31,95,168}
\definecolor{PlotRed}{RGB}{200,55,45}
\definecolor{PlotGreen}{RGB}{40,140,70}
\definecolor{PlotOrange}{RGB}{225,130,20}
\definecolor{PlotPurple}{RGB}{120,60,160}
\definecolor{PlotGray}{RGB}{120,120,120}

%----------------------------------------------------------------------
% Graphics and plotting
%----------------------------------------------------------------------
\usepackage{graphicx}
\usepackage{tikz}
\usetikzlibrary{arrows.meta,calc,positioning,decorations.pathreplacing,
  patterns,shapes.geometric,angles,quotes,backgrounds,fit,matrix}
\usepackage{pgfplots}
\pgfplotsset{compat=1.18}
\usepgfplotslibrary{fillbetween,groupplots}
% A common look for all function plots in the notes.
\pgfplotsset{
  mlplot/.style={
    axis lines=middle,
    axis line style={-{Stealth[length=2mm]},thin},
    xlabel style={anchor=north west,font=\small},
    ylabel style={anchor=south east,font=\small},
    tick label style={font=\footnotesize},
    every axis plot/.append style={thick},
    grid=major,
    grid style={gray!20},
    clip=true,
  },
}
\tikzset{
  vec/.style={-{Stealth[length=2.6mm]},very thick},
  dot/.style={circle,fill,inner sep=1.3pt},
  opendot/.style={circle,draw,fill=white,inner sep=1.3pt},
}

%----------------------------------------------------------------------
% Tables, lists, captions, code
%----------------------------------------------------------------------
\usepackage{booktabs}
\usepackage{array}
\usepackage{tabularx}
\usepackage{multirow}
\usepackage{enumitem}
\setlist{itemsep=2pt,topsep=3pt,parsep=0pt}
\setlist[enumerate,1]{label=(\alph*)}
\newlist{steps}{enumerate}{1}
\setlist[steps]{label=\textbf{Step \arabic*.},leftmargin=*,align=left,
  itemsep=3pt,topsep=3pt}
\usepackage[font=small,labelfont={bf,color=MLchapter},margin=8mm]{caption}
\usepackage{subcaption}
\usepackage{listings}
\lstdefinestyle{mlpython}{
  language=Python,
  basicstyle=\ttfamily\small,
  keywordstyle=\color{MLdef}\bfseries,
  commentstyle=\color{MLremark}\itshape,
  stringstyle=\color{MLexample},
  showstringspaces=false,
  columns=fullflexible,
  keepspaces=true,
  frame=single,
  framesep=4pt,
  rulecolor=\color{gray!40},
  backgroundcolor=\color{gray!5},
  xleftmargin=6pt,
  xrightmargin=6pt,
  aboveskip=8pt,
  belowskip=8pt,
  upquote=true,
}
\lstset{style=mlpython}

%----------------------------------------------------------------------
% Headings and running heads
%----------------------------------------------------------------------
\usepackage{titlesec}
% The last command of the "before" code receives the chapter title;
% \MLsetchaptertitle prints it and remembers it for the solutions
% appendix (see the exercise machinery below).
\newcommand{\MLcurrentchaptertitle}{}
\newcommand{\MLsetchaptertitle}[1]{%
  \gdef\MLcurrentchaptertitle{#1}#1}
\titleformat{\chapter}[display]
  {\normalfont\color{MLchapter}\filright}
  {\large\scshape\chaptertitlename\ \thechapter}
  {4pt}
  {\huge\bfseries\MLsetchaptertitle}
  [\vspace{2pt}{\color{MLchapter}\titlerule[0.8pt]}]
\titleformat{name=\chapter,numberless}[display]
  {\normalfont\color{MLchapter}\filright}
  {}
  {0pt}
  {\huge\bfseries}
  [\vspace{2pt}{\color{MLchapter}\titlerule[0.8pt]}]
\titlespacing*{\chapter}{0pt}{-10pt}{22pt}
\titleformat{\section}
  {\normalfont\Large\bfseries\color{MLchapter}\filright}
  {\thesection}{0.7em}{}
\titleformat{\subsection}
  {\normalfont\large\bfseries\color{MLchapter}\filright}
  {\thesubsection}{0.7em}{}
\titleformat{\subsubsection}
  {\normalfont\normalsize\bfseries\filright}
  {\thesubsubsection}{0.7em}{}
\titlespacing*{\section}{0pt}{16pt plus 4pt minus 2pt}{6pt plus 2pt}
\titlespacing*{\subsection}{0pt}{12pt plus 3pt minus 2pt}{4pt plus 2pt}
\titlespacing*{\subsubsection}{0pt}{10pt plus 3pt minus 2pt}{3pt plus 1pt}
\setcounter{secnumdepth}{2}
\setcounter{tocdepth}{1}

\usepackage{fancyhdr}
\pagestyle{fancy}
\fancyhf{}
\fancyhead[L]{\small\itshape\nouppercase{\leftmark}}
\fancyhead[R]{\small\thepage}
\renewcommand{\headrulewidth}{0.4pt}
\fancypagestyle{plain}{%
  \fancyhf{}%
  \fancyfoot[C]{\small\thepage}%
  \renewcommand{\headrulewidth}{0pt}}
% Running head: "4. Vectors and Linear Algebra ..." (no number in the
% front matter). Must come after \pagestyle{fancy}, which resets it.
\makeatletter
\renewcommand{\chaptermark}[1]{%
  \markboth{\ifnum\c@secnumdepth>\m@ne\if@mainmatter
    \thechapter.\enspace\fi\fi #1}{}}
\makeatother

%----------------------------------------------------------------------
% Index, hyperlinks, clever references
%----------------------------------------------------------------------
\usepackage{etoolbox}
\usepackage{imakeidx}
\makeindex[intoc,columns=2,title=Index]
\usepackage[hyperfootnotes=false]{hyperref}
\hypersetup{
  colorlinks=true,
  linkcolor=MLlink,
  citecolor=MLlink,
  urlcolor=MLurl,
  bookmarksnumbered=true,
  bookmarksopen=true,
  bookmarksopenlevel=0,
  pdftitle={Mathematics for Machine Learning and Deep Learning: From
    Absolute Beginner to Prerequisites},
  pdfsubject={Lecture notes},
  pdfkeywords={mathematics, machine learning, deep learning, linear algebra,
    calculus, probability, statistics, optimization},
}
\usepackage[capitalise,nameinlink,noabbrev]{cleveref}

%----------------------------------------------------------------------
% Boxes
%----------------------------------------------------------------------
\usepackage{tcolorbox}
\tcbuselibrary{skins,breakable,theorems}

% Base style shared by every coloured box: a thick rule on the left, a
% lightly tinted background, and the ability to break across pages.
\tcbset{
  mlbox/.style={
    enhanced,
    breakable,
    sharp corners,
    boxrule=0pt,
    frame hidden,
    colback=#1!6!white,
    borderline west={2.5pt}{0pt}{#1},
    left=9pt,right=8pt,top=6pt,bottom=6pt,
    before skip=9pt plus 2pt,
    after skip=9pt plus 2pt,
    parbox=false,
  },
}

%----------------------------------------------------------------------
% Theorem-like environments
%
% Numbered (one shared counter per chapter, so "Example 4.7" follows
% "Definition 4.6"):  definition, theorem, proposition, corollary,
%                     example, algorithm
% Numbered separately: exercise
% Unnumbered:         intuition, picturethis, remark, mlnote, pitfall,
%                     motivation, notation
%
% Every one of them is drawn inside a tcolorbox (see below), and every
% one accepts an optional title, e.g. \begin{definition}[Slope].
%----------------------------------------------------------------------
\newcommand{\MLnewthmstyle}[2]{%
  \newtheoremstyle{#1}%
    {0pt}{0pt}%                 space above / below (the box handles it)
    {\normalfont}%              body font
    {}%                         indentation of the head
    {\bfseries\color{#2}}%      head font
    {.}%                        punctuation after the head
    {0.5em}%                    space after the head
    {}}%                        head specification (default)
\MLnewthmstyle{mldefinition}{MLdef}
\MLnewthmstyle{mltheorem}{MLthm}
\MLnewthmstyle{mlexample}{MLexample}
\MLnewthmstyle{mlintuition}{MLintuition}
\MLnewthmstyle{mlml}{MLml}
\MLnewthmstyle{mlpitfall}{MLpitfall}
\MLnewthmstyle{mlremark}{MLremark}
\MLnewthmstyle{mlmotivation}{MLmotivation}
\MLnewthmstyle{mlnotation}{MLnotation}
\MLnewthmstyle{mlexercise}{MLexercise}
\MLnewthmstyle{mlalgorithm}{MLalgorithm}

\theoremstyle{mldefinition}
\newtheorem{definition}{Definition}[chapter]
\theoremstyle{mltheorem}
\newtheorem{theorem}[definition]{Theorem}
\newtheorem{proposition}[definition]{Proposition}
\newtheorem{corollary}[definition]{Corollary}
\theoremstyle{mlexample}
\newtheorem{example}[definition]{Example}
\theoremstyle{mlalgorithm}
\newtheorem{algorithm}[definition]{Algorithm}
\theoremstyle{mlintuition}
\newtheorem*{intuition}{Intuition}
\newtheorem*{picturethis}{Picture this}
\theoremstyle{mlml}
\newtheorem*{mlnote}{Why this matters for ML}
\theoremstyle{mlpitfall}
\newtheorem*{pitfall}{Common pitfall}
\theoremstyle{mlremark}
\newtheorem*{remark}{Remark}
\theoremstyle{mlmotivation}
\newtheorem*{motivation}{Motivation}
\theoremstyle{mlnotation}
\newtheorem*{notation}{Notation}
\theoremstyle{mlexercise}
\newtheorem{exercisebase}{Exercise}[chapter]

\crefname{definition}{Definition}{Definitions}
\crefname{theorem}{Theorem}{Theorems}
\crefname{proposition}{Proposition}{Propositions}
\crefname{corollary}{Corollary}{Corollaries}
\crefname{example}{Example}{Examples}
\crefname{algorithm}{Algorithm}{Algorithms}
\crefname{exercisebase}{Exercise}{Exercises}
\crefname{figure}{Figure}{Figures}
\crefname{table}{Table}{Tables}
\crefname{equation}{Equation}{Equations}
\crefname{chapter}{Chapter}{Chapters}
\crefname{section}{Section}{Sections}
\crefname{subsection}{Section}{Sections}
\crefname{appendix}{Appendix}{Appendices}

\tcolorboxenvironment{definition}{mlbox=MLdef}
\tcolorboxenvironment{theorem}{mlbox=MLthm}
\tcolorboxenvironment{proposition}{mlbox=MLthm}
\tcolorboxenvironment{corollary}{mlbox=MLthm}
\tcolorboxenvironment{example}{mlbox=MLexample,colback=MLexample!3!white}
\tcolorboxenvironment{algorithm}{mlbox=MLalgorithm}
\tcolorboxenvironment{intuition}{mlbox=MLintuition,colback=MLintuition!7!white}
\tcolorboxenvironment{picturethis}{mlbox=MLintuition,colback=MLintuition!7!white}
\tcolorboxenvironment{mlnote}{mlbox=MLml}
\tcolorboxenvironment{pitfall}{mlbox=MLpitfall}
\tcolorboxenvironment{remark}{mlbox=MLremark,colback=MLremark!4!white}
\tcolorboxenvironment{motivation}{mlbox=MLmotivation}
\tcolorboxenvironment{notation}{mlbox=MLnotation}

% Proofs: amsthm's proof environment with a coloured QED box.
\renewcommand{\qedsymbol}{\textcolor{MLthm}{$\blacksquare$}}

% "Solution." label used inside worked examples.
\newcommand{\sol}{\par\smallskip\noindent\textit{Solution.}\enspace}

% A box for the learning goals at the start of each chapter.
\newtcolorbox{objectives}{
  enhanced,breakable,sharp corners,
  colback=MLchapter!4!white,colframe=MLchapter!60!white,boxrule=0.6pt,
  title={In this chapter you will learn to\dots},
  fonttitle=\bfseries\color{white},colbacktitle=MLchapter!85!white,
  left=8pt,right=8pt,top=5pt,bottom=5pt,before skip=10pt,after skip=12pt,
  parbox=false,
}

% A box for the summary at the end of each chapter.
\newtcolorbox{keytakeaways}[1][Key takeaways]{
  enhanced,breakable,sharp corners,
  colback=MLsummary!4!white,colframe=MLsummary!60!white,boxrule=0.6pt,
  title={#1},
  fonttitle=\bfseries\color{white},colbacktitle=MLsummary!85!white,
  left=8pt,right=8pt,top=5pt,bottom=5pt,before skip=10pt,after skip=10pt,
  parbox=false,
}

% A frame around a key formula, for use inside display maths:
%   \[ \keyeq{a^2+b^2=c^2} \]
% (The options are passed explicitly because tcolorbox resets the global
% "highlight math style" at the start of every box.)
\tcbset{mlkeyeq/.style={enhanced,sharp corners,colback=white,
  colframe=MLchapter!70!white,boxrule=0.8pt,left=5pt,right=5pt,
  top=3pt,bottom=3pt}}
\newcommand{\keyeq}[1]{\tcbhighmath[mlkeyeq]{#1}}

%----------------------------------------------------------------------
% New terms: bold, coloured, and automatically added to the index.
%   \term{gradient}               -> index entry "gradient"
%   \term[mean!arithmetic]{mean}  -> custom index entry
%----------------------------------------------------------------------
\newcommand{\term}[2][]{%
  \textbf{\color{MLterm}#2}%
  \ifblank{#1}{\index{#2}}{\index{#1}}}

%----------------------------------------------------------------------
% Exercises and automatically collected solutions
%
% Usage inside a chapter:
%
%   \sectionexercises
%   \begin{exercise}
%     Compute $3+4\times 2$.
%   \end{exercise}
%   \begin{solution}
%     Multiplication first: $3+8=11$.
%   \end{solution}
%
% The solution is not printed where it is written. Instead it is stored
% in the file \jobname.sol, together with the number of the exercise it
% belongs to, and \printsolutions (used in the solutions appendix)
% typesets all of them, grouped by chapter, with links in both
% directions. Solutions must not contain verbatim material or the
% character #.
%----------------------------------------------------------------------
\makeatletter
\newcounter{MLsolutioncount}
\newwrite\ML@solfile
\newif\ifML@solfileopen
\gdef\ML@lastsolchapter{}

\AtBeginDocument{%
  \immediate\openout\ML@solfile=\jobname.sol\relax
  \global\ML@solfileopentrue}

\newenvironment{exercise}[1][]{%
  \ifblank{#1}{\begin{exercisebase}}{\begin{exercisebase}[#1]}%
  \label{exercise:\theexercisebase}}%
  {\end{exercisebase}}
% The link to the solution sits in the bottom-right corner, in a strip
% of extra padding so that it never overlaps the text of the exercise.
\tcolorboxenvironment{exercise}{%
  mlbox=MLexercise,colback=MLexercise!3!white,bottom=12pt,
  overlay unbroken and last={%
    \node[anchor=south east,inner sep=3pt,font=\scriptsize]
      at (frame.south east)
      {\hyperref[solution:\theexercisebase]{\textcolor{MLexercise!75!white}%
        {Solution on p.\,\pageref*{solution:\theexercisebase}\,%
          $\rightarrow$}}};}}

\NewDocumentEnvironment{solution}{+b}{%
  \ifML@solfileopen
    \edef\ML@thischapter{\thechapter}%
    \ifx\ML@thischapter\ML@lastsolchapter\else
      \immediate\write\ML@solfile{%
        \string\MLSolutionsChapter{\thechapter}%
        {\unexpanded\expandafter{\MLcurrentchaptertitle}}}%
      \global\let\ML@lastsolchapter\ML@thischapter
    \fi
    \immediate\write\ML@solfile{%
      \string\MLSolution{\theexercisebase}{\unexpanded{#1}}}%
    \stepcounter{MLsolutioncount}%
  \fi}{}

% Heading printed before the exercises of a section.
\newcommand{\sectionexercises}{%
  \par\needspace{6\baselineskip}\medskip
  \noindent{\large\bfseries\color{MLexercise}Exercises for
    Section~\thesection}\par\nopagebreak\smallskip}

% Typesetting of the stored solutions.
\newcommand{\MLSolutionsChapter}[2]{%
  \section*{Chapter #1: #2}%
  \phantomsection
  \addcontentsline{toc}{section}{Chapter #1: #2}}
\newcommand{\MLSolution}[2]{%
  \par\needspace{3\baselineskip}\medskip\noindent
  \phantomsection\def\@currentlabel{#1}\label{solution:#1}%
  {\bfseries\hyperref[exercise:#1]{\textcolor{MLexercise}{Exercise~#1.}}}%
  \enspace #2\par}
\newcommand{\printsolutions}{%
  \ifML@solfileopen
    \immediate\closeout\ML@solfile
    \global\ML@solfileopenfalse
  \fi
  \InputIfFileExists{\jobname.sol}{}{%
    \par\textit{No solutions were recorded in this run.}\par}}
\makeatother

%----------------------------------------------------------------------
% Notation macros (see Appendix A for the full list)
%----------------------------------------------------------------------
% Number sets
\newcommand{\N}{\mathbb{N}}
\newcommand{\Z}{\mathbb{Z}}
\newcommand{\Q}{\mathbb{Q}}
\newcommand{\R}{\mathbb{R}}

% Vectors (bold lower case) and matrices (bold upper case)
\newcommand{\vect}[1]{\bm{#1}}
\newcommand{\mat}[1]{\bm{#1}}
\newcommand{\va}{\vect{a}}
\newcommand{\vb}{\vect{b}}
\newcommand{\vc}{\vect{c}}
\newcommand{\vd}{\vect{d}}
\newcommand{\ve}{\vect{e}}
\newcommand{\vg}{\vect{g}}
\newcommand{\vh}{\vect{h}}
\newcommand{\vk}{\vect{k}}
\newcommand{\vm}{\vect{m}}
\newcommand{\vp}{\vect{p}}
\newcommand{\vq}{\vect{q}}
\newcommand{\vr}{\vect{r}}
\newcommand{\vs}{\vect{s}}
\newcommand{\vu}{\vect{u}}
\newcommand{\vv}{\vect{v}}
\newcommand{\vw}{\vect{w}}
\newcommand{\vx}{\vect{x}}
\newcommand{\vy}{\vect{y}}
\newcommand{\vz}{\vect{z}}
\newcommand{\vzero}{\vect{0}}
\newcommand{\vone}{\vect{1}}
\newcommand{\vtheta}{\bm{\theta}}
\newcommand{\vmu}{\bm{\mu}}
\newcommand{\vdelta}{\bm{\delta}}
\newcommand{\vepsilon}{\bm{\varepsilon}}
\newcommand{\mA}{\mat{A}}
\newcommand{\mB}{\mat{B}}
\newcommand{\mC}{\mat{C}}
\newcommand{\mD}{\mat{D}}
\newcommand{\mH}{\mat{H}}
\newcommand{\mI}{\mat{I}}
\newcommand{\mM}{\mat{M}}
\newcommand{\mP}{\mat{P}}
\newcommand{\mQ}{\mat{Q}}
\newcommand{\mR}{\mat{R}}
\newcommand{\mS}{\mat{S}}
\newcommand{\mU}{\mat{U}}
\newcommand{\mV}{\mat{V}}
\newcommand{\mW}{\mat{W}}
\newcommand{\mX}{\mat{X}}
\newcommand{\mY}{\mat{Y}}
\newcommand{\mZero}{\mat{0}}
\newcommand{\mSigma}{\bm{\Sigma}}
\newcommand{\mLambda}{\bm{\Lambda}}

% Delimiters
\DeclarePairedDelimiter{\abs}{\lvert}{\rvert}
\DeclarePairedDelimiter{\norm}{\lVert}{\rVert}
\DeclarePairedDelimiter{\ceil}{\lceil}{\rceil}
\DeclarePairedDelimiter{\floor}{\lfloor}{\rfloor}

% Operators
\DeclareMathOperator*{\argmin}{arg\,min}
\DeclareMathOperator*{\argmax}{arg\,max}
\DeclareMathOperator{\Var}{Var}
\DeclareMathOperator{\Cov}{Cov}
\DeclareMathOperator{\Corr}{Corr}
\DeclareMathOperator{\Bias}{Bias}
\DeclareMathOperator{\MSE}{MSE}
\DeclareMathOperator{\tr}{tr}
\DeclareMathOperator{\rank}{rank}
\DeclareMathOperator{\diag}{diag}
\DeclareMathOperator{\sign}{sign}
\DeclareMathOperator{\spn}{span}
\DeclareMathOperator{\proj}{proj}
\DeclareMathOperator{\softmax}{softmax}
\DeclareMathOperator{\ReLU}{ReLU}
\DeclareMathOperator{\KL}{KL}
\newcommand{\E}{\mathbb{E}}
\newcommand{\dd}{\mathrm{d}}
\newcommand{\deriv}[2]{\frac{\dd #1}{\dd #2}}
\newcommand{\pderiv}[2]{\frac{\partial #1}{\partial #2}}
\newcommand{\iid}{\overset{\text{i.i.d.}}{\sim}}
. Everything that controls the
% look of the notes (fonts, colours, boxes, headings), the notation
% macros, and the automatic exercise-solution machinery lives here, so
% that every chapter file contains only content.
%======================================================================

%----------------------------------------------------------------------
% Fonts, encoding and micro-typography
%----------------------------------------------------------------------
\usepackage[T1]{fontenc}
\usepackage{lmodern}
\usepackage[final]{microtype}

%----------------------------------------------------------------------
% Page geometry and paragraph layout
%----------------------------------------------------------------------
\usepackage[a4paper,top=27mm,bottom=27mm,left=30mm,right=30mm,
            headheight=14pt]{geometry}
\usepackage[skip=5pt plus 1pt minus 1pt,indent=0pt]{parskip}
\usepackage{needspace}

%----------------------------------------------------------------------
% Mathematics
%----------------------------------------------------------------------
\usepackage{amsmath}
\usepackage{amssymb}
\usepackage{amsthm}
\usepackage{mathtools}
\usepackage{bm}

%----------------------------------------------------------------------
% Colours (one colour per kind of box, used consistently throughout)
%----------------------------------------------------------------------
\usepackage{xcolor}
\definecolor{MLchapter}{RGB}{25,60,110}      % chapter and section titles
\definecolor{MLlink}{RGB}{20,80,160}         % internal hyperlinks
\definecolor{MLurl}{RGB}{140,40,110}         % external hyperlinks
\definecolor{MLterm}{RGB}{20,55,120}         % newly introduced terms
\definecolor{MLdef}{RGB}{31,95,168}          % definitions
\definecolor{MLthm}{RGB}{128,40,110}         % theorems, propositions
\definecolor{MLexample}{RGB}{40,120,60}      % worked examples
\definecolor{MLintuition}{RGB}{176,104,0}    % intuition / picture this
\definecolor{MLml}{RGB}{0,118,124}           % why this matters for ML
\definecolor{MLpitfall}{RGB}{185,45,40}      % common pitfalls
\definecolor{MLremark}{RGB}{100,100,100}     % remarks
\definecolor{MLmotivation}{RGB}{70,60,160}   % motivation
\definecolor{MLnotation}{RGB}{70,100,115}    % notation
\definecolor{MLexercise}{RGB}{60,60,60}      % exercises
\definecolor{MLalgorithm}{RGB}{120,80,40}    % algorithms
\definecolor{MLsummary}{RGB}{25,60,110}      % chapter summaries
% Colours used inside figures
\definecolor{PlotBlue}{RGB}{31,95,168}
\definecolor{PlotRed}{RGB}{200,55,45}
\definecolor{PlotGreen}{RGB}{40,140,70}
\definecolor{PlotOrange}{RGB}{225,130,20}
\definecolor{PlotPurple}{RGB}{120,60,160}
\definecolor{PlotGray}{RGB}{120,120,120}

%----------------------------------------------------------------------
% Graphics and plotting
%----------------------------------------------------------------------
\usepackage{graphicx}
\usepackage{tikz}
\usetikzlibrary{arrows.meta,calc,positioning,decorations.pathreplacing,
  patterns,shapes.geometric,angles,quotes,backgrounds,fit,matrix}
\usepackage{pgfplots}
\pgfplotsset{compat=1.18}
\usepgfplotslibrary{fillbetween,groupplots}
% A common look for all function plots in the notes.
\pgfplotsset{
  mlplot/.style={
    axis lines=middle,
    axis line style={-{Stealth[length=2mm]},thin},
    xlabel style={anchor=north west,font=\small},
    ylabel style={anchor=south east,font=\small},
    tick label style={font=\footnotesize},
    every axis plot/.append style={thick},
    grid=major,
    grid style={gray!20},
    clip=true,
  },
}
\tikzset{
  vec/.style={-{Stealth[length=2.6mm]},very thick},
  dot/.style={circle,fill,inner sep=1.3pt},
  opendot/.style={circle,draw,fill=white,inner sep=1.3pt},
}

%----------------------------------------------------------------------
% Tables, lists, captions, code
%----------------------------------------------------------------------
\usepackage{booktabs}
\usepackage{array}
\usepackage{tabularx}
\usepackage{multirow}
\usepackage{enumitem}
\setlist{itemsep=2pt,topsep=3pt,parsep=0pt}
\setlist[enumerate,1]{label=(\alph*)}
\newlist{steps}{enumerate}{1}
\setlist[steps]{label=\textbf{Step \arabic*.},leftmargin=*,align=left,
  itemsep=3pt,topsep=3pt}
\usepackage[font=small,labelfont={bf,color=MLchapter},margin=8mm]{caption}
\usepackage{subcaption}
\usepackage{listings}
\lstdefinestyle{mlpython}{
  language=Python,
  basicstyle=\ttfamily\small,
  keywordstyle=\color{MLdef}\bfseries,
  commentstyle=\color{MLremark}\itshape,
  stringstyle=\color{MLexample},
  showstringspaces=false,
  columns=fullflexible,
  keepspaces=true,
  frame=single,
  framesep=4pt,
  rulecolor=\color{gray!40},
  backgroundcolor=\color{gray!5},
  xleftmargin=6pt,
  xrightmargin=6pt,
  aboveskip=8pt,
  belowskip=8pt,
  upquote=true,
}
\lstset{style=mlpython}

%----------------------------------------------------------------------
% Headings and running heads
%----------------------------------------------------------------------
\usepackage{titlesec}
% The last command of the "before" code receives the chapter title;
% \MLsetchaptertitle prints it and remembers it for the solutions
% appendix (see the exercise machinery below).
\newcommand{\MLcurrentchaptertitle}{}
\newcommand{\MLsetchaptertitle}[1]{%
  \gdef\MLcurrentchaptertitle{#1}#1}
\titleformat{\chapter}[display]
  {\normalfont\color{MLchapter}\filright}
  {\large\scshape\chaptertitlename\ \thechapter}
  {4pt}
  {\huge\bfseries\MLsetchaptertitle}
  [\vspace{2pt}{\color{MLchapter}\titlerule[0.8pt]}]
\titleformat{name=\chapter,numberless}[display]
  {\normalfont\color{MLchapter}\filright}
  {}
  {0pt}
  {\huge\bfseries}
  [\vspace{2pt}{\color{MLchapter}\titlerule[0.8pt]}]
\titlespacing*{\chapter}{0pt}{-10pt}{22pt}
\titleformat{\section}
  {\normalfont\Large\bfseries\color{MLchapter}\filright}
  {\thesection}{0.7em}{}
\titleformat{\subsection}
  {\normalfont\large\bfseries\color{MLchapter}\filright}
  {\thesubsection}{0.7em}{}
\titleformat{\subsubsection}
  {\normalfont\normalsize\bfseries\filright}
  {\thesubsubsection}{0.7em}{}
\titlespacing*{\section}{0pt}{16pt plus 4pt minus 2pt}{6pt plus 2pt}
\titlespacing*{\subsection}{0pt}{12pt plus 3pt minus 2pt}{4pt plus 2pt}
\titlespacing*{\subsubsection}{0pt}{10pt plus 3pt minus 2pt}{3pt plus 1pt}
\setcounter{secnumdepth}{2}
\setcounter{tocdepth}{1}

\usepackage{fancyhdr}
\pagestyle{fancy}
\fancyhf{}
\fancyhead[L]{\small\itshape\nouppercase{\leftmark}}
\fancyhead[R]{\small\thepage}
\renewcommand{\headrulewidth}{0.4pt}
\fancypagestyle{plain}{%
  \fancyhf{}%
  \fancyfoot[C]{\small\thepage}%
  \renewcommand{\headrulewidth}{0pt}}
% Running head: "4. Vectors and Linear Algebra ..." (no number in the
% front matter). Must come after \pagestyle{fancy}, which resets it.
\makeatletter
\renewcommand{\chaptermark}[1]{%
  \markboth{\ifnum\c@secnumdepth>\m@ne\if@mainmatter
    \thechapter.\enspace\fi\fi #1}{}}
\makeatother

%----------------------------------------------------------------------
% Index, hyperlinks, clever references
%----------------------------------------------------------------------
\usepackage{etoolbox}
\usepackage{imakeidx}
\makeindex[intoc,columns=2,title=Index]
\usepackage[hyperfootnotes=false]{hyperref}
\hypersetup{
  colorlinks=true,
  linkcolor=MLlink,
  citecolor=MLlink,
  urlcolor=MLurl,
  bookmarksnumbered=true,
  bookmarksopen=true,
  bookmarksopenlevel=0,
  pdftitle={Mathematics for Machine Learning and Deep Learning: From
    Absolute Beginner to Prerequisites},
  pdfsubject={Lecture notes},
  pdfkeywords={mathematics, machine learning, deep learning, linear algebra,
    calculus, probability, statistics, optimization},
}
\usepackage[capitalise,nameinlink,noabbrev]{cleveref}

%----------------------------------------------------------------------
% Boxes
%----------------------------------------------------------------------
\usepackage{tcolorbox}
\tcbuselibrary{skins,breakable,theorems}

% Base style shared by every coloured box: a thick rule on the left, a
% lightly tinted background, and the ability to break across pages.
\tcbset{
  mlbox/.style={
    enhanced,
    breakable,
    sharp corners,
    boxrule=0pt,
    frame hidden,
    colback=#1!6!white,
    borderline west={2.5pt}{0pt}{#1},
    left=9pt,right=8pt,top=6pt,bottom=6pt,
    before skip=9pt plus 2pt,
    after skip=9pt plus 2pt,
    parbox=false,
  },
}

%----------------------------------------------------------------------
% Theorem-like environments
%
% Numbered (one shared counter per chapter, so "Example 4.7" follows
% "Definition 4.6"):  definition, theorem, proposition, corollary,
%                     example, algorithm
% Numbered separately: exercise
% Unnumbered:         intuition, picturethis, remark, mlnote, pitfall,
%                     motivation, notation
%
% Every one of them is drawn inside a tcolorbox (see below), and every
% one accepts an optional title, e.g. \begin{definition}[Slope].
%----------------------------------------------------------------------
\newcommand{\MLnewthmstyle}[2]{%
  \newtheoremstyle{#1}%
    {0pt}{0pt}%                 space above / below (the box handles it)
    {\normalfont}%              body font
    {}%                         indentation of the head
    {\bfseries\color{#2}}%      head font
    {.}%                        punctuation after the head
    {0.5em}%                    space after the head
    {}}%                        head specification (default)
\MLnewthmstyle{mldefinition}{MLdef}
\MLnewthmstyle{mltheorem}{MLthm}
\MLnewthmstyle{mlexample}{MLexample}
\MLnewthmstyle{mlintuition}{MLintuition}
\MLnewthmstyle{mlml}{MLml}
\MLnewthmstyle{mlpitfall}{MLpitfall}
\MLnewthmstyle{mlremark}{MLremark}
\MLnewthmstyle{mlmotivation}{MLmotivation}
\MLnewthmstyle{mlnotation}{MLnotation}
\MLnewthmstyle{mlexercise}{MLexercise}
\MLnewthmstyle{mlalgorithm}{MLalgorithm}

\theoremstyle{mldefinition}
\newtheorem{definition}{Definition}[chapter]
\theoremstyle{mltheorem}
\newtheorem{theorem}[definition]{Theorem}
\newtheorem{proposition}[definition]{Proposition}
\newtheorem{corollary}[definition]{Corollary}
\theoremstyle{mlexample}
\newtheorem{example}[definition]{Example}
\theoremstyle{mlalgorithm}
\newtheorem{algorithm}[definition]{Algorithm}
\theoremstyle{mlintuition}
\newtheorem*{intuition}{Intuition}
\newtheorem*{picturethis}{Picture this}
\theoremstyle{mlml}
\newtheorem*{mlnote}{Why this matters for ML}
\theoremstyle{mlpitfall}
\newtheorem*{pitfall}{Common pitfall}
\theoremstyle{mlremark}
\newtheorem*{remark}{Remark}
\theoremstyle{mlmotivation}
\newtheorem*{motivation}{Motivation}
\theoremstyle{mlnotation}
\newtheorem*{notation}{Notation}
\theoremstyle{mlexercise}
\newtheorem{exercisebase}{Exercise}[chapter]

\crefname{definition}{Definition}{Definitions}
\crefname{theorem}{Theorem}{Theorems}
\crefname{proposition}{Proposition}{Propositions}
\crefname{corollary}{Corollary}{Corollaries}
\crefname{example}{Example}{Examples}
\crefname{algorithm}{Algorithm}{Algorithms}
\crefname{exercisebase}{Exercise}{Exercises}
\crefname{figure}{Figure}{Figures}
\crefname{table}{Table}{Tables}
\crefname{equation}{Equation}{Equations}
\crefname{chapter}{Chapter}{Chapters}
\crefname{section}{Section}{Sections}
\crefname{subsection}{Section}{Sections}
\crefname{appendix}{Appendix}{Appendices}

\tcolorboxenvironment{definition}{mlbox=MLdef}
\tcolorboxenvironment{theorem}{mlbox=MLthm}
\tcolorboxenvironment{proposition}{mlbox=MLthm}
\tcolorboxenvironment{corollary}{mlbox=MLthm}
\tcolorboxenvironment{example}{mlbox=MLexample,colback=MLexample!3!white}
\tcolorboxenvironment{algorithm}{mlbox=MLalgorithm}
\tcolorboxenvironment{intuition}{mlbox=MLintuition,colback=MLintuition!7!white}
\tcolorboxenvironment{picturethis}{mlbox=MLintuition,colback=MLintuition!7!white}
\tcolorboxenvironment{mlnote}{mlbox=MLml}
\tcolorboxenvironment{pitfall}{mlbox=MLpitfall}
\tcolorboxenvironment{remark}{mlbox=MLremark,colback=MLremark!4!white}
\tcolorboxenvironment{motivation}{mlbox=MLmotivation}
\tcolorboxenvironment{notation}{mlbox=MLnotation}

% Proofs: amsthm's proof environment with a coloured QED box.
\renewcommand{\qedsymbol}{\textcolor{MLthm}{$\blacksquare$}}

% "Solution." label used inside worked examples.
\newcommand{\sol}{\par\smallskip\noindent\textit{Solution.}\enspace}

% A box for the learning goals at the start of each chapter.
\newtcolorbox{objectives}{
  enhanced,breakable,sharp corners,
  colback=MLchapter!4!white,colframe=MLchapter!60!white,boxrule=0.6pt,
  title={In this chapter you will learn to\dots},
  fonttitle=\bfseries\color{white},colbacktitle=MLchapter!85!white,
  left=8pt,right=8pt,top=5pt,bottom=5pt,before skip=10pt,after skip=12pt,
  parbox=false,
}

% A box for the summary at the end of each chapter.
\newtcolorbox{keytakeaways}[1][Key takeaways]{
  enhanced,breakable,sharp corners,
  colback=MLsummary!4!white,colframe=MLsummary!60!white,boxrule=0.6pt,
  title={#1},
  fonttitle=\bfseries\color{white},colbacktitle=MLsummary!85!white,
  left=8pt,right=8pt,top=5pt,bottom=5pt,before skip=10pt,after skip=10pt,
  parbox=false,
}

% A frame around a key formula, for use inside display maths:
%   \[ \keyeq{a^2+b^2=c^2} \]
% (The options are passed explicitly because tcolorbox resets the global
% "highlight math style" at the start of every box.)
\tcbset{mlkeyeq/.style={enhanced,sharp corners,colback=white,
  colframe=MLchapter!70!white,boxrule=0.8pt,left=5pt,right=5pt,
  top=3pt,bottom=3pt}}
\newcommand{\keyeq}[1]{\tcbhighmath[mlkeyeq]{#1}}

%----------------------------------------------------------------------
% New terms: bold, coloured, and automatically added to the index.
%   \term{gradient}               -> index entry "gradient"
%   \term[mean!arithmetic]{mean}  -> custom index entry
%----------------------------------------------------------------------
\newcommand{\term}[2][]{%
  \textbf{\color{MLterm}#2}%
  \ifblank{#1}{\index{#2}}{\index{#1}}}

%----------------------------------------------------------------------
% Exercises and automatically collected solutions
%
% Usage inside a chapter:
%
%   \sectionexercises
%   \begin{exercise}
%     Compute $3+4\times 2$.
%   \end{exercise}
%   \begin{solution}
%     Multiplication first: $3+8=11$.
%   \end{solution}
%
% The solution is not printed where it is written. Instead it is stored
% in the file \jobname.sol, together with the number of the exercise it
% belongs to, and \printsolutions (used in the solutions appendix)
% typesets all of them, grouped by chapter, with links in both
% directions. Solutions must not contain verbatim material or the
% character #.
%----------------------------------------------------------------------
\makeatletter
\newcounter{MLsolutioncount}
\newwrite\ML@solfile
\newif\ifML@solfileopen
\gdef\ML@lastsolchapter{}

\AtBeginDocument{%
  \immediate\openout\ML@solfile=\jobname.sol\relax
  \global\ML@solfileopentrue}

\newenvironment{exercise}[1][]{%
  \ifblank{#1}{\begin{exercisebase}}{\begin{exercisebase}[#1]}%
  \label{exercise:\theexercisebase}}%
  {\end{exercisebase}}
% The link to the solution sits in the bottom-right corner, in a strip
% of extra padding so that it never overlaps the text of the exercise.
\tcolorboxenvironment{exercise}{%
  mlbox=MLexercise,colback=MLexercise!3!white,bottom=12pt,
  overlay unbroken and last={%
    \node[anchor=south east,inner sep=3pt,font=\scriptsize]
      at (frame.south east)
      {\hyperref[solution:\theexercisebase]{\textcolor{MLexercise!75!white}%
        {Solution on p.\,\pageref*{solution:\theexercisebase}\,%
          $\rightarrow$}}};}}

\NewDocumentEnvironment{solution}{+b}{%
  \ifML@solfileopen
    \edef\ML@thischapter{\thechapter}%
    \ifx\ML@thischapter\ML@lastsolchapter\else
      \immediate\write\ML@solfile{%
        \string\MLSolutionsChapter{\thechapter}%
        {\unexpanded\expandafter{\MLcurrentchaptertitle}}}%
      \global\let\ML@lastsolchapter\ML@thischapter
    \fi
    \immediate\write\ML@solfile{%
      \string\MLSolution{\theexercisebase}{\unexpanded{#1}}}%
    \stepcounter{MLsolutioncount}%
  \fi}{}

% Heading printed before the exercises of a section.
\newcommand{\sectionexercises}{%
  \par\needspace{6\baselineskip}\medskip
  \noindent{\large\bfseries\color{MLexercise}Exercises for
    Section~\thesection}\par\nopagebreak\smallskip}

% Typesetting of the stored solutions.
\newcommand{\MLSolutionsChapter}[2]{%
  \section*{Chapter #1: #2}%
  \phantomsection
  \addcontentsline{toc}{section}{Chapter #1: #2}}
\newcommand{\MLSolution}[2]{%
  \par\needspace{3\baselineskip}\medskip\noindent
  \phantomsection\def\@currentlabel{#1}\label{solution:#1}%
  {\bfseries\hyperref[exercise:#1]{\textcolor{MLexercise}{Exercise~#1.}}}%
  \enspace #2\par}
\newcommand{\printsolutions}{%
  \ifML@solfileopen
    \immediate\closeout\ML@solfile
    \global\ML@solfileopenfalse
  \fi
  \InputIfFileExists{\jobname.sol}{}{%
    \par\textit{No solutions were recorded in this run.}\par}}
\makeatother

%----------------------------------------------------------------------
% Notation macros (see Appendix A for the full list)
%----------------------------------------------------------------------
% Number sets
\newcommand{\N}{\mathbb{N}}
\newcommand{\Z}{\mathbb{Z}}
\newcommand{\Q}{\mathbb{Q}}
\newcommand{\R}{\mathbb{R}}

% Vectors (bold lower case) and matrices (bold upper case)
\newcommand{\vect}[1]{\bm{#1}}
\newcommand{\mat}[1]{\bm{#1}}
\newcommand{\va}{\vect{a}}
\newcommand{\vb}{\vect{b}}
\newcommand{\vc}{\vect{c}}
\newcommand{\vd}{\vect{d}}
\newcommand{\ve}{\vect{e}}
\newcommand{\vg}{\vect{g}}
\newcommand{\vh}{\vect{h}}
\newcommand{\vk}{\vect{k}}
\newcommand{\vm}{\vect{m}}
\newcommand{\vp}{\vect{p}}
\newcommand{\vq}{\vect{q}}
\newcommand{\vr}{\vect{r}}
\newcommand{\vs}{\vect{s}}
\newcommand{\vu}{\vect{u}}
\newcommand{\vv}{\vect{v}}
\newcommand{\vw}{\vect{w}}
\newcommand{\vx}{\vect{x}}
\newcommand{\vy}{\vect{y}}
\newcommand{\vz}{\vect{z}}
\newcommand{\vzero}{\vect{0}}
\newcommand{\vone}{\vect{1}}
\newcommand{\vtheta}{\bm{\theta}}
\newcommand{\vmu}{\bm{\mu}}
\newcommand{\vdelta}{\bm{\delta}}
\newcommand{\vepsilon}{\bm{\varepsilon}}
\newcommand{\mA}{\mat{A}}
\newcommand{\mB}{\mat{B}}
\newcommand{\mC}{\mat{C}}
\newcommand{\mD}{\mat{D}}
\newcommand{\mH}{\mat{H}}
\newcommand{\mI}{\mat{I}}
\newcommand{\mM}{\mat{M}}
\newcommand{\mP}{\mat{P}}
\newcommand{\mQ}{\mat{Q}}
\newcommand{\mR}{\mat{R}}
\newcommand{\mS}{\mat{S}}
\newcommand{\mU}{\mat{U}}
\newcommand{\mV}{\mat{V}}
\newcommand{\mW}{\mat{W}}
\newcommand{\mX}{\mat{X}}
\newcommand{\mY}{\mat{Y}}
\newcommand{\mZero}{\mat{0}}
\newcommand{\mSigma}{\bm{\Sigma}}
\newcommand{\mLambda}{\bm{\Lambda}}

% Delimiters
\DeclarePairedDelimiter{\abs}{\lvert}{\rvert}
\DeclarePairedDelimiter{\norm}{\lVert}{\rVert}
\DeclarePairedDelimiter{\ceil}{\lceil}{\rceil}
\DeclarePairedDelimiter{\floor}{\lfloor}{\rfloor}

% Operators
\DeclareMathOperator*{\argmin}{arg\,min}
\DeclareMathOperator*{\argmax}{arg\,max}
\DeclareMathOperator{\Var}{Var}
\DeclareMathOperator{\Cov}{Cov}
\DeclareMathOperator{\Corr}{Corr}
\DeclareMathOperator{\Bias}{Bias}
\DeclareMathOperator{\MSE}{MSE}
\DeclareMathOperator{\tr}{tr}
\DeclareMathOperator{\rank}{rank}
\DeclareMathOperator{\diag}{diag}
\DeclareMathOperator{\sign}{sign}
\DeclareMathOperator{\spn}{span}
\DeclareMathOperator{\proj}{proj}
\DeclareMathOperator{\softmax}{softmax}
\DeclareMathOperator{\ReLU}{ReLU}
\DeclareMathOperator{\KL}{KL}
\newcommand{\E}{\mathbb{E}}
\newcommand{\dd}{\mathrm{d}}
\newcommand{\deriv}[2]{\frac{\dd #1}{\dd #2}}
\newcommand{\pderiv}[2]{\frac{\partial #1}{\partial #2}}
\newcommand{\iid}{\overset{\text{i.i.d.}}{\sim}}
. Everything that controls the
% look of the notes (fonts, colours, boxes, headings), the notation
% macros, and the automatic exercise-solution machinery lives here, so
% that every chapter file contains only content.
%======================================================================

%----------------------------------------------------------------------
% Fonts, encoding and micro-typography
%----------------------------------------------------------------------
\usepackage[T1]{fontenc}
\usepackage{lmodern}
\usepackage[final]{microtype}

%----------------------------------------------------------------------
% Page geometry and paragraph layout
%----------------------------------------------------------------------
\usepackage[a4paper,top=27mm,bottom=27mm,left=30mm,right=30mm,
            headheight=14pt]{geometry}
\usepackage[skip=5pt plus 1pt minus 1pt,indent=0pt]{parskip}
\usepackage{needspace}

%----------------------------------------------------------------------
% Mathematics
%----------------------------------------------------------------------
\usepackage{amsmath}
\usepackage{amssymb}
\usepackage{amsthm}
\usepackage{mathtools}
\usepackage{bm}

%----------------------------------------------------------------------
% Colours (one colour per kind of box, used consistently throughout)
%----------------------------------------------------------------------
\usepackage{xcolor}
\definecolor{MLchapter}{RGB}{25,60,110}      % chapter and section titles
\definecolor{MLlink}{RGB}{20,80,160}         % internal hyperlinks
\definecolor{MLurl}{RGB}{140,40,110}         % external hyperlinks
\definecolor{MLterm}{RGB}{20,55,120}         % newly introduced terms
\definecolor{MLdef}{RGB}{31,95,168}          % definitions
\definecolor{MLthm}{RGB}{128,40,110}         % theorems, propositions
\definecolor{MLexample}{RGB}{40,120,60}      % worked examples
\definecolor{MLintuition}{RGB}{176,104,0}    % intuition / picture this
\definecolor{MLml}{RGB}{0,118,124}           % why this matters for ML
\definecolor{MLpitfall}{RGB}{185,45,40}      % common pitfalls
\definecolor{MLremark}{RGB}{100,100,100}     % remarks
\definecolor{MLmotivation}{RGB}{70,60,160}   % motivation
\definecolor{MLnotation}{RGB}{70,100,115}    % notation
\definecolor{MLexercise}{RGB}{60,60,60}      % exercises
\definecolor{MLalgorithm}{RGB}{120,80,40}    % algorithms
\definecolor{MLsummary}{RGB}{25,60,110}      % chapter summaries
% Colours used inside figures
\definecolor{PlotBlue}{RGB}{31,95,168}
\definecolor{PlotRed}{RGB}{200,55,45}
\definecolor{PlotGreen}{RGB}{40,140,70}
\definecolor{PlotOrange}{RGB}{225,130,20}
\definecolor{PlotPurple}{RGB}{120,60,160}
\definecolor{PlotGray}{RGB}{120,120,120}

%----------------------------------------------------------------------
% Graphics and plotting
%----------------------------------------------------------------------
\usepackage{graphicx}
\usepackage{tikz}
\usetikzlibrary{arrows.meta,calc,positioning,decorations.pathreplacing,
  patterns,shapes.geometric,angles,quotes,backgrounds,fit,matrix}
\usepackage{pgfplots}
\pgfplotsset{compat=1.18}
\usepgfplotslibrary{fillbetween,groupplots}
% A common look for all function plots in the notes.
\pgfplotsset{
  mlplot/.style={
    axis lines=middle,
    axis line style={-{Stealth[length=2mm]},thin},
    xlabel style={anchor=north west,font=\small},
    ylabel style={anchor=south east,font=\small},
    tick label style={font=\footnotesize},
    every axis plot/.append style={thick},
    grid=major,
    grid style={gray!20},
    clip=true,
  },
}
\tikzset{
  vec/.style={-{Stealth[length=2.6mm]},very thick},
  dot/.style={circle,fill,inner sep=1.3pt},
  opendot/.style={circle,draw,fill=white,inner sep=1.3pt},
}

%----------------------------------------------------------------------
% Tables, lists, captions, code
%----------------------------------------------------------------------
\usepackage{booktabs}
\usepackage{array}
\usepackage{tabularx}
\usepackage{multirow}
\usepackage{enumitem}
\setlist{itemsep=2pt,topsep=3pt,parsep=0pt}
\setlist[enumerate,1]{label=(\alph*)}
\newlist{steps}{enumerate}{1}
\setlist[steps]{label=\textbf{Step \arabic*.},leftmargin=*,align=left,
  itemsep=3pt,topsep=3pt}
\usepackage[font=small,labelfont={bf,color=MLchapter},margin=8mm]{caption}
\usepackage{subcaption}
\usepackage{listings}
\lstdefinestyle{mlpython}{
  language=Python,
  basicstyle=\ttfamily\small,
  keywordstyle=\color{MLdef}\bfseries,
  commentstyle=\color{MLremark}\itshape,
  stringstyle=\color{MLexample},
  showstringspaces=false,
  columns=fullflexible,
  keepspaces=true,
  frame=single,
  framesep=4pt,
  rulecolor=\color{gray!40},
  backgroundcolor=\color{gray!5},
  xleftmargin=6pt,
  xrightmargin=6pt,
  aboveskip=8pt,
  belowskip=8pt,
  upquote=true,
}
\lstset{style=mlpython}

%----------------------------------------------------------------------
% Headings and running heads
%----------------------------------------------------------------------
\usepackage{titlesec}
% The last command of the "before" code receives the chapter title;
% \MLsetchaptertitle prints it and remembers it for the solutions
% appendix (see the exercise machinery below).
\newcommand{\MLcurrentchaptertitle}{}
\newcommand{\MLsetchaptertitle}[1]{%
  \gdef\MLcurrentchaptertitle{#1}#1}
\titleformat{\chapter}[display]
  {\normalfont\color{MLchapter}\filright}
  {\large\scshape\chaptertitlename\ \thechapter}
  {4pt}
  {\huge\bfseries\MLsetchaptertitle}
  [\vspace{2pt}{\color{MLchapter}\titlerule[0.8pt]}]
\titleformat{name=\chapter,numberless}[display]
  {\normalfont\color{MLchapter}\filright}
  {}
  {0pt}
  {\huge\bfseries}
  [\vspace{2pt}{\color{MLchapter}\titlerule[0.8pt]}]
\titlespacing*{\chapter}{0pt}{-10pt}{22pt}
\titleformat{\section}
  {\normalfont\Large\bfseries\color{MLchapter}\filright}
  {\thesection}{0.7em}{}
\titleformat{\subsection}
  {\normalfont\large\bfseries\color{MLchapter}\filright}
  {\thesubsection}{0.7em}{}
\titleformat{\subsubsection}
  {\normalfont\normalsize\bfseries\filright}
  {\thesubsubsection}{0.7em}{}
\titlespacing*{\section}{0pt}{16pt plus 4pt minus 2pt}{6pt plus 2pt}
\titlespacing*{\subsection}{0pt}{12pt plus 3pt minus 2pt}{4pt plus 2pt}
\titlespacing*{\subsubsection}{0pt}{10pt plus 3pt minus 2pt}{3pt plus 1pt}
\setcounter{secnumdepth}{2}
\setcounter{tocdepth}{1}

\usepackage{fancyhdr}
\pagestyle{fancy}
\fancyhf{}
\fancyhead[L]{\small\itshape\nouppercase{\leftmark}}
\fancyhead[R]{\small\thepage}
\renewcommand{\headrulewidth}{0.4pt}
\fancypagestyle{plain}{%
  \fancyhf{}%
  \fancyfoot[C]{\small\thepage}%
  \renewcommand{\headrulewidth}{0pt}}
% Running head: "4. Vectors and Linear Algebra ..." (no number in the
% front matter). Must come after \pagestyle{fancy}, which resets it.
\makeatletter
\renewcommand{\chaptermark}[1]{%
  \markboth{\ifnum\c@secnumdepth>\m@ne\if@mainmatter
    \thechapter.\enspace\fi\fi #1}{}}
\makeatother

%----------------------------------------------------------------------
% Index, hyperlinks, clever references
%----------------------------------------------------------------------
\usepackage{etoolbox}
\usepackage{imakeidx}
\makeindex[intoc,columns=2,title=Index]
\usepackage[hyperfootnotes=false]{hyperref}
\hypersetup{
  colorlinks=true,
  linkcolor=MLlink,
  citecolor=MLlink,
  urlcolor=MLurl,
  bookmarksnumbered=true,
  bookmarksopen=true,
  bookmarksopenlevel=0,
  pdftitle={Mathematics for Machine Learning and Deep Learning: From
    Absolute Beginner to Prerequisites},
  pdfsubject={Lecture notes},
  pdfkeywords={mathematics, machine learning, deep learning, linear algebra,
    calculus, probability, statistics, optimization},
}
\usepackage[capitalise,nameinlink,noabbrev]{cleveref}

%----------------------------------------------------------------------
% Boxes
%----------------------------------------------------------------------
\usepackage{tcolorbox}
\tcbuselibrary{skins,breakable,theorems}

% Base style shared by every coloured box: a thick rule on the left, a
% lightly tinted background, and the ability to break across pages.
\tcbset{
  mlbox/.style={
    enhanced,
    breakable,
    sharp corners,
    boxrule=0pt,
    frame hidden,
    colback=#1!6!white,
    borderline west={2.5pt}{0pt}{#1},
    left=9pt,right=8pt,top=6pt,bottom=6pt,
    before skip=9pt plus 2pt,
    after skip=9pt plus 2pt,
    parbox=false,
  },
}

%----------------------------------------------------------------------
% Theorem-like environments
%
% Numbered (one shared counter per chapter, so "Example 4.7" follows
% "Definition 4.6"):  definition, theorem, proposition, corollary,
%                     example, algorithm
% Numbered separately: exercise
% Unnumbered:         intuition, picturethis, remark, mlnote, pitfall,
%                     motivation, notation
%
% Every one of them is drawn inside a tcolorbox (see below), and every
% one accepts an optional title, e.g. \begin{definition}[Slope].
%----------------------------------------------------------------------
\newcommand{\MLnewthmstyle}[2]{%
  \newtheoremstyle{#1}%
    {0pt}{0pt}%                 space above / below (the box handles it)
    {\normalfont}%              body font
    {}%                         indentation of the head
    {\bfseries\color{#2}}%      head font
    {.}%                        punctuation after the head
    {0.5em}%                    space after the head
    {}}%                        head specification (default)
\MLnewthmstyle{mldefinition}{MLdef}
\MLnewthmstyle{mltheorem}{MLthm}
\MLnewthmstyle{mlexample}{MLexample}
\MLnewthmstyle{mlintuition}{MLintuition}
\MLnewthmstyle{mlml}{MLml}
\MLnewthmstyle{mlpitfall}{MLpitfall}
\MLnewthmstyle{mlremark}{MLremark}
\MLnewthmstyle{mlmotivation}{MLmotivation}
\MLnewthmstyle{mlnotation}{MLnotation}
\MLnewthmstyle{mlexercise}{MLexercise}
\MLnewthmstyle{mlalgorithm}{MLalgorithm}

\theoremstyle{mldefinition}
\newtheorem{definition}{Definition}[chapter]
\theoremstyle{mltheorem}
\newtheorem{theorem}[definition]{Theorem}
\newtheorem{proposition}[definition]{Proposition}
\newtheorem{corollary}[definition]{Corollary}
\theoremstyle{mlexample}
\newtheorem{example}[definition]{Example}
\theoremstyle{mlalgorithm}
\newtheorem{algorithm}[definition]{Algorithm}
\theoremstyle{mlintuition}
\newtheorem*{intuition}{Intuition}
\newtheorem*{picturethis}{Picture this}
\theoremstyle{mlml}
\newtheorem*{mlnote}{Why this matters for ML}
\theoremstyle{mlpitfall}
\newtheorem*{pitfall}{Common pitfall}
\theoremstyle{mlremark}
\newtheorem*{remark}{Remark}
\theoremstyle{mlmotivation}
\newtheorem*{motivation}{Motivation}
\theoremstyle{mlnotation}
\newtheorem*{notation}{Notation}
\theoremstyle{mlexercise}
\newtheorem{exercisebase}{Exercise}[chapter]

\crefname{definition}{Definition}{Definitions}
\crefname{theorem}{Theorem}{Theorems}
\crefname{proposition}{Proposition}{Propositions}
\crefname{corollary}{Corollary}{Corollaries}
\crefname{example}{Example}{Examples}
\crefname{algorithm}{Algorithm}{Algorithms}
\crefname{exercisebase}{Exercise}{Exercises}
\crefname{figure}{Figure}{Figures}
\crefname{table}{Table}{Tables}
\crefname{equation}{Equation}{Equations}
\crefname{chapter}{Chapter}{Chapters}
\crefname{section}{Section}{Sections}
\crefname{subsection}{Section}{Sections}
\crefname{appendix}{Appendix}{Appendices}

\tcolorboxenvironment{definition}{mlbox=MLdef}
\tcolorboxenvironment{theorem}{mlbox=MLthm}
\tcolorboxenvironment{proposition}{mlbox=MLthm}
\tcolorboxenvironment{corollary}{mlbox=MLthm}
\tcolorboxenvironment{example}{mlbox=MLexample,colback=MLexample!3!white}
\tcolorboxenvironment{algorithm}{mlbox=MLalgorithm}
\tcolorboxenvironment{intuition}{mlbox=MLintuition,colback=MLintuition!7!white}
\tcolorboxenvironment{picturethis}{mlbox=MLintuition,colback=MLintuition!7!white}
\tcolorboxenvironment{mlnote}{mlbox=MLml}
\tcolorboxenvironment{pitfall}{mlbox=MLpitfall}
\tcolorboxenvironment{remark}{mlbox=MLremark,colback=MLremark!4!white}
\tcolorboxenvironment{motivation}{mlbox=MLmotivation}
\tcolorboxenvironment{notation}{mlbox=MLnotation}

% Proofs: amsthm's proof environment with a coloured QED box.
\renewcommand{\qedsymbol}{\textcolor{MLthm}{$\blacksquare$}}

% "Solution." label used inside worked examples.
\newcommand{\sol}{\par\smallskip\noindent\textit{Solution.}\enspace}

% A box for the learning goals at the start of each chapter.
\newtcolorbox{objectives}{
  enhanced,breakable,sharp corners,
  colback=MLchapter!4!white,colframe=MLchapter!60!white,boxrule=0.6pt,
  title={In this chapter you will learn to\dots},
  fonttitle=\bfseries\color{white},colbacktitle=MLchapter!85!white,
  left=8pt,right=8pt,top=5pt,bottom=5pt,before skip=10pt,after skip=12pt,
  parbox=false,
}

% A box for the summary at the end of each chapter.
\newtcolorbox{keytakeaways}[1][Key takeaways]{
  enhanced,breakable,sharp corners,
  colback=MLsummary!4!white,colframe=MLsummary!60!white,boxrule=0.6pt,
  title={#1},
  fonttitle=\bfseries\color{white},colbacktitle=MLsummary!85!white,
  left=8pt,right=8pt,top=5pt,bottom=5pt,before skip=10pt,after skip=10pt,
  parbox=false,
}

% A frame around a key formula, for use inside display maths:
%   \[ \keyeq{a^2+b^2=c^2} \]
% (The options are passed explicitly because tcolorbox resets the global
% "highlight math style" at the start of every box.)
\tcbset{mlkeyeq/.style={enhanced,sharp corners,colback=white,
  colframe=MLchapter!70!white,boxrule=0.8pt,left=5pt,right=5pt,
  top=3pt,bottom=3pt}}
\newcommand{\keyeq}[1]{\tcbhighmath[mlkeyeq]{#1}}

%----------------------------------------------------------------------
% New terms: bold, coloured, and automatically added to the index.
%   \term{gradient}               -> index entry "gradient"
%   \term[mean!arithmetic]{mean}  -> custom index entry
%----------------------------------------------------------------------
\newcommand{\term}[2][]{%
  \textbf{\color{MLterm}#2}%
  \ifblank{#1}{\index{#2}}{\index{#1}}}

%----------------------------------------------------------------------
% Exercises and automatically collected solutions
%
% Usage inside a chapter:
%
%   \sectionexercises
%   \begin{exercise}
%     Compute $3+4\times 2$.
%   \end{exercise}
%   \begin{solution}
%     Multiplication first: $3+8=11$.
%   \end{solution}
%
% The solution is not printed where it is written. Instead it is stored
% in the file \jobname.sol, together with the number of the exercise it
% belongs to, and \printsolutions (used in the solutions appendix)
% typesets all of them, grouped by chapter, with links in both
% directions. Solutions must not contain verbatim material or the
% character #.
%----------------------------------------------------------------------
\makeatletter
\newcounter{MLsolutioncount}
\newwrite\ML@solfile
\newif\ifML@solfileopen
\gdef\ML@lastsolchapter{}

\AtBeginDocument{%
  \immediate\openout\ML@solfile=\jobname.sol\relax
  \global\ML@solfileopentrue}

\newenvironment{exercise}[1][]{%
  \ifblank{#1}{\begin{exercisebase}}{\begin{exercisebase}[#1]}%
  \label{exercise:\theexercisebase}}%
  {\end{exercisebase}}
% The link to the solution sits in the bottom-right corner, in a strip
% of extra padding so that it never overlaps the text of the exercise.
\tcolorboxenvironment{exercise}{%
  mlbox=MLexercise,colback=MLexercise!3!white,bottom=12pt,
  overlay unbroken and last={%
    \node[anchor=south east,inner sep=3pt,font=\scriptsize]
      at (frame.south east)
      {\hyperref[solution:\theexercisebase]{\textcolor{MLexercise!75!white}%
        {Solution on p.\,\pageref*{solution:\theexercisebase}\,%
          $\rightarrow$}}};}}

\NewDocumentEnvironment{solution}{+b}{%
  \ifML@solfileopen
    \edef\ML@thischapter{\thechapter}%
    \ifx\ML@thischapter\ML@lastsolchapter\else
      \immediate\write\ML@solfile{%
        \string\MLSolutionsChapter{\thechapter}%
        {\unexpanded\expandafter{\MLcurrentchaptertitle}}}%
      \global\let\ML@lastsolchapter\ML@thischapter
    \fi
    \immediate\write\ML@solfile{%
      \string\MLSolution{\theexercisebase}{\unexpanded{#1}}}%
    \stepcounter{MLsolutioncount}%
  \fi}{}

% Heading printed before the exercises of a section.
\newcommand{\sectionexercises}{%
  \par\needspace{6\baselineskip}\medskip
  \noindent{\large\bfseries\color{MLexercise}Exercises for
    Section~\thesection}\par\nopagebreak\smallskip}

% Typesetting of the stored solutions.
\newcommand{\MLSolutionsChapter}[2]{%
  \section*{Chapter #1: #2}%
  \phantomsection
  \addcontentsline{toc}{section}{Chapter #1: #2}}
\newcommand{\MLSolution}[2]{%
  \par\needspace{3\baselineskip}\medskip\noindent
  \phantomsection\def\@currentlabel{#1}\label{solution:#1}%
  {\bfseries\hyperref[exercise:#1]{\textcolor{MLexercise}{Exercise~#1.}}}%
  \enspace #2\par}
\newcommand{\printsolutions}{%
  \ifML@solfileopen
    \immediate\closeout\ML@solfile
    \global\ML@solfileopenfalse
  \fi
  \InputIfFileExists{\jobname.sol}{}{%
    \par\textit{No solutions were recorded in this run.}\par}}
\makeatother

%----------------------------------------------------------------------
% Notation macros (see Appendix A for the full list)
%----------------------------------------------------------------------
% Number sets
\newcommand{\N}{\mathbb{N}}
\newcommand{\Z}{\mathbb{Z}}
\newcommand{\Q}{\mathbb{Q}}
\newcommand{\R}{\mathbb{R}}

% Vectors (bold lower case) and matrices (bold upper case)
\newcommand{\vect}[1]{\bm{#1}}
\newcommand{\mat}[1]{\bm{#1}}
\newcommand{\va}{\vect{a}}
\newcommand{\vb}{\vect{b}}
\newcommand{\vc}{\vect{c}}
\newcommand{\vd}{\vect{d}}
\newcommand{\ve}{\vect{e}}
\newcommand{\vg}{\vect{g}}
\newcommand{\vh}{\vect{h}}
\newcommand{\vk}{\vect{k}}
\newcommand{\vm}{\vect{m}}
\newcommand{\vp}{\vect{p}}
\newcommand{\vq}{\vect{q}}
\newcommand{\vr}{\vect{r}}
\newcommand{\vs}{\vect{s}}
\newcommand{\vu}{\vect{u}}
\newcommand{\vv}{\vect{v}}
\newcommand{\vw}{\vect{w}}
\newcommand{\vx}{\vect{x}}
\newcommand{\vy}{\vect{y}}
\newcommand{\vz}{\vect{z}}
\newcommand{\vzero}{\vect{0}}
\newcommand{\vone}{\vect{1}}
\newcommand{\vtheta}{\bm{\theta}}
\newcommand{\vmu}{\bm{\mu}}
\newcommand{\vdelta}{\bm{\delta}}
\newcommand{\vepsilon}{\bm{\varepsilon}}
\newcommand{\mA}{\mat{A}}
\newcommand{\mB}{\mat{B}}
\newcommand{\mC}{\mat{C}}
\newcommand{\mD}{\mat{D}}
\newcommand{\mH}{\mat{H}}
\newcommand{\mI}{\mat{I}}
\newcommand{\mM}{\mat{M}}
\newcommand{\mP}{\mat{P}}
\newcommand{\mQ}{\mat{Q}}
\newcommand{\mR}{\mat{R}}
\newcommand{\mS}{\mat{S}}
\newcommand{\mU}{\mat{U}}
\newcommand{\mV}{\mat{V}}
\newcommand{\mW}{\mat{W}}
\newcommand{\mX}{\mat{X}}
\newcommand{\mY}{\mat{Y}}
\newcommand{\mZero}{\mat{0}}
\newcommand{\mSigma}{\bm{\Sigma}}
\newcommand{\mLambda}{\bm{\Lambda}}

% Delimiters
\DeclarePairedDelimiter{\abs}{\lvert}{\rvert}
\DeclarePairedDelimiter{\norm}{\lVert}{\rVert}
\DeclarePairedDelimiter{\ceil}{\lceil}{\rceil}
\DeclarePairedDelimiter{\floor}{\lfloor}{\rfloor}

% Operators
\DeclareMathOperator*{\argmin}{arg\,min}
\DeclareMathOperator*{\argmax}{arg\,max}
\DeclareMathOperator{\Var}{Var}
\DeclareMathOperator{\Cov}{Cov}
\DeclareMathOperator{\Corr}{Corr}
\DeclareMathOperator{\Bias}{Bias}
\DeclareMathOperator{\MSE}{MSE}
\DeclareMathOperator{\tr}{tr}
\DeclareMathOperator{\rank}{rank}
\DeclareMathOperator{\diag}{diag}
\DeclareMathOperator{\sign}{sign}
\DeclareMathOperator{\spn}{span}
\DeclareMathOperator{\proj}{proj}
\DeclareMathOperator{\softmax}{softmax}
\DeclareMathOperator{\ReLU}{ReLU}
\DeclareMathOperator{\KL}{KL}
\newcommand{\E}{\mathbb{E}}
\newcommand{\dd}{\mathrm{d}}
\newcommand{\deriv}[2]{\frac{\dd #1}{\dd #2}}
\newcommand{\pderiv}[2]{\frac{\partial #1}{\partial #2}}
\newcommand{\iid}{\overset{\text{i.i.d.}}{\sim}}
. Everything that controls the
% look of the notes (fonts, colours, boxes, headings), the notation
% macros, and the automatic exercise-solution machinery lives here, so
% that every chapter file contains only content.
%======================================================================

%----------------------------------------------------------------------
% Fonts, encoding and micro-typography
%----------------------------------------------------------------------
\usepackage[T1]{fontenc}
\usepackage{lmodern}
\usepackage[final]{microtype}

%----------------------------------------------------------------------
% Page geometry and paragraph layout
%----------------------------------------------------------------------
\usepackage[a4paper,top=27mm,bottom=27mm,left=30mm,right=30mm,
            headheight=14pt]{geometry}
\usepackage[skip=5pt plus 1pt minus 1pt,indent=0pt]{parskip}
\usepackage{needspace}

%----------------------------------------------------------------------
% Mathematics
%----------------------------------------------------------------------
\usepackage{amsmath}
\usepackage{amssymb}
\usepackage{amsthm}
\usepackage{mathtools}
\usepackage{bm}

%----------------------------------------------------------------------
% Colours (one colour per kind of box, used consistently throughout)
%----------------------------------------------------------------------
\usepackage{xcolor}
\definecolor{MLchapter}{RGB}{25,60,110}      % chapter and section titles
\definecolor{MLlink}{RGB}{20,80,160}         % internal hyperlinks
\definecolor{MLurl}{RGB}{140,40,110}         % external hyperlinks
\definecolor{MLterm}{RGB}{20,55,120}         % newly introduced terms
\definecolor{MLdef}{RGB}{31,95,168}          % definitions
\definecolor{MLthm}{RGB}{128,40,110}         % theorems, propositions
\definecolor{MLexample}{RGB}{40,120,60}      % worked examples
\definecolor{MLintuition}{RGB}{176,104,0}    % intuition / picture this
\definecolor{MLml}{RGB}{0,118,124}           % why this matters for ML
\definecolor{MLpitfall}{RGB}{185,45,40}      % common pitfalls
\definecolor{MLremark}{RGB}{100,100,100}     % remarks
\definecolor{MLmotivation}{RGB}{70,60,160}   % motivation
\definecolor{MLnotation}{RGB}{70,100,115}    % notation
\definecolor{MLexercise}{RGB}{60,60,60}      % exercises
\definecolor{MLalgorithm}{RGB}{120,80,40}    % algorithms
\definecolor{MLsummary}{RGB}{25,60,110}      % chapter summaries
% Colours used inside figures
\definecolor{PlotBlue}{RGB}{31,95,168}
\definecolor{PlotRed}{RGB}{200,55,45}
\definecolor{PlotGreen}{RGB}{40,140,70}
\definecolor{PlotOrange}{RGB}{225,130,20}
\definecolor{PlotPurple}{RGB}{120,60,160}
\definecolor{PlotGray}{RGB}{120,120,120}

%----------------------------------------------------------------------
% Graphics and plotting
%----------------------------------------------------------------------
\usepackage{graphicx}
\usepackage{tikz}
\usetikzlibrary{arrows.meta,calc,positioning,decorations.pathreplacing,
  patterns,shapes.geometric,angles,quotes,backgrounds,fit,matrix}
\usepackage{pgfplots}
\pgfplotsset{compat=1.18}
\usepgfplotslibrary{fillbetween,groupplots}
% A common look for all function plots in the notes.
\pgfplotsset{
  mlplot/.style={
    axis lines=middle,
    axis line style={-{Stealth[length=2mm]},thin},
    xlabel style={anchor=north west,font=\small},
    ylabel style={anchor=south east,font=\small},
    tick label style={font=\footnotesize},
    every axis plot/.append style={thick},
    grid=major,
    grid style={gray!20},
    clip=true,
  },
}
\tikzset{
  vec/.style={-{Stealth[length=2.6mm]},very thick},
  dot/.style={circle,fill,inner sep=1.3pt},
  opendot/.style={circle,draw,fill=white,inner sep=1.3pt},
}

%----------------------------------------------------------------------
% Tables, lists, captions, code
%----------------------------------------------------------------------
\usepackage{booktabs}
\usepackage{array}
\usepackage{tabularx}
\usepackage{multirow}
\usepackage{enumitem}
\setlist{itemsep=2pt,topsep=3pt,parsep=0pt}
\setlist[enumerate,1]{label=(\alph*)}
\newlist{steps}{enumerate}{1}
\setlist[steps]{label=\textbf{Step \arabic*.},leftmargin=*,align=left,
  itemsep=3pt,topsep=3pt}
\usepackage[font=small,labelfont={bf,color=MLchapter},margin=8mm]{caption}
\usepackage{subcaption}
\usepackage{listings}
\lstdefinestyle{mlpython}{
  language=Python,
  basicstyle=\ttfamily\small,
  keywordstyle=\color{MLdef}\bfseries,
  commentstyle=\color{MLremark}\itshape,
  stringstyle=\color{MLexample},
  showstringspaces=false,
  columns=fullflexible,
  keepspaces=true,
  frame=single,
  framesep=4pt,
  rulecolor=\color{gray!40},
  backgroundcolor=\color{gray!5},
  xleftmargin=6pt,
  xrightmargin=6pt,
  aboveskip=8pt,
  belowskip=8pt,
  upquote=true,
}
\lstset{style=mlpython}

%----------------------------------------------------------------------
% Headings and running heads
%----------------------------------------------------------------------
\usepackage{titlesec}
% The last command of the "before" code receives the chapter title;
% \MLsetchaptertitle prints it and remembers it for the solutions
% appendix (see the exercise machinery below).
\newcommand{\MLcurrentchaptertitle}{}
\newcommand{\MLsetchaptertitle}[1]{%
  \gdef\MLcurrentchaptertitle{#1}#1}
\titleformat{\chapter}[display]
  {\normalfont\color{MLchapter}\filright}
  {\large\scshape\chaptertitlename\ \thechapter}
  {4pt}
  {\huge\bfseries\MLsetchaptertitle}
  [\vspace{2pt}{\color{MLchapter}\titlerule[0.8pt]}]
\titleformat{name=\chapter,numberless}[display]
  {\normalfont\color{MLchapter}\filright}
  {}
  {0pt}
  {\huge\bfseries}
  [\vspace{2pt}{\color{MLchapter}\titlerule[0.8pt]}]
\titlespacing*{\chapter}{0pt}{-10pt}{22pt}
\titleformat{\section}
  {\normalfont\Large\bfseries\color{MLchapter}\filright}
  {\thesection}{0.7em}{}
\titleformat{\subsection}
  {\normalfont\large\bfseries\color{MLchapter}\filright}
  {\thesubsection}{0.7em}{}
\titleformat{\subsubsection}
  {\normalfont\normalsize\bfseries\filright}
  {\thesubsubsection}{0.7em}{}
\titlespacing*{\section}{0pt}{16pt plus 4pt minus 2pt}{6pt plus 2pt}
\titlespacing*{\subsection}{0pt}{12pt plus 3pt minus 2pt}{4pt plus 2pt}
\titlespacing*{\subsubsection}{0pt}{10pt plus 3pt minus 2pt}{3pt plus 1pt}
\setcounter{secnumdepth}{2}
\setcounter{tocdepth}{1}

\usepackage{fancyhdr}
\pagestyle{fancy}
\fancyhf{}
\fancyhead[L]{\small\itshape\nouppercase{\leftmark}}
\fancyhead[R]{\small\thepage}
\renewcommand{\headrulewidth}{0.4pt}
\fancypagestyle{plain}{%
  \fancyhf{}%
  \fancyfoot[C]{\small\thepage}%
  \renewcommand{\headrulewidth}{0pt}}
% Running head: "4. Vectors and Linear Algebra ..." (no number in the
% front matter). Must come after \pagestyle{fancy}, which resets it.
\makeatletter
\renewcommand{\chaptermark}[1]{%
  \markboth{\ifnum\c@secnumdepth>\m@ne\if@mainmatter
    \thechapter.\enspace\fi\fi #1}{}}
\makeatother

%----------------------------------------------------------------------
% Index, hyperlinks, clever references
%----------------------------------------------------------------------
\usepackage{etoolbox}
\usepackage{imakeidx}
\makeindex[intoc,columns=2,title=Index]
\usepackage[hyperfootnotes=false]{hyperref}
\hypersetup{
  colorlinks=true,
  linkcolor=MLlink,
  citecolor=MLlink,
  urlcolor=MLurl,
  bookmarksnumbered=true,
  bookmarksopen=true,
  bookmarksopenlevel=0,
  pdftitle={Mathematics for Machine Learning and Deep Learning: From
    Absolute Beginner to Prerequisites},
  pdfsubject={Lecture notes},
  pdfkeywords={mathematics, machine learning, deep learning, linear algebra,
    calculus, probability, statistics, optimization},
}
\usepackage[capitalise,nameinlink,noabbrev]{cleveref}

%----------------------------------------------------------------------
% Boxes
%----------------------------------------------------------------------
\usepackage{tcolorbox}
\tcbuselibrary{skins,breakable,theorems}

% Base style shared by every coloured box: a thick rule on the left, a
% lightly tinted background, and the ability to break across pages.
\tcbset{
  mlbox/.style={
    enhanced,
    breakable,
    sharp corners,
    boxrule=0pt,
    frame hidden,
    colback=#1!6!white,
    borderline west={2.5pt}{0pt}{#1},
    left=9pt,right=8pt,top=6pt,bottom=6pt,
    before skip=9pt plus 2pt,
    after skip=9pt plus 2pt,
    parbox=false,
  },
}

%----------------------------------------------------------------------
% Theorem-like environments
%
% Numbered (one shared counter per chapter, so "Example 4.7" follows
% "Definition 4.6"):  definition, theorem, proposition, corollary,
%                     example, algorithm
% Numbered separately: exercise
% Unnumbered:         intuition, picturethis, remark, mlnote, pitfall,
%                     motivation, notation
%
% Every one of them is drawn inside a tcolorbox (see below), and every
% one accepts an optional title, e.g. \begin{definition}[Slope].
%----------------------------------------------------------------------
\newcommand{\MLnewthmstyle}[2]{%
  \newtheoremstyle{#1}%
    {0pt}{0pt}%                 space above / below (the box handles it)
    {\normalfont}%              body font
    {}%                         indentation of the head
    {\bfseries\color{#2}}%      head font
    {.}%                        punctuation after the head
    {0.5em}%                    space after the head
    {}}%                        head specification (default)
\MLnewthmstyle{mldefinition}{MLdef}
\MLnewthmstyle{mltheorem}{MLthm}
\MLnewthmstyle{mlexample}{MLexample}
\MLnewthmstyle{mlintuition}{MLintuition}
\MLnewthmstyle{mlml}{MLml}
\MLnewthmstyle{mlpitfall}{MLpitfall}
\MLnewthmstyle{mlremark}{MLremark}
\MLnewthmstyle{mlmotivation}{MLmotivation}
\MLnewthmstyle{mlnotation}{MLnotation}
\MLnewthmstyle{mlexercise}{MLexercise}
\MLnewthmstyle{mlalgorithm}{MLalgorithm}

\theoremstyle{mldefinition}
\newtheorem{definition}{Definition}[chapter]
\theoremstyle{mltheorem}
\newtheorem{theorem}[definition]{Theorem}
\newtheorem{proposition}[definition]{Proposition}
\newtheorem{corollary}[definition]{Corollary}
\theoremstyle{mlexample}
\newtheorem{example}[definition]{Example}
\theoremstyle{mlalgorithm}
\newtheorem{algorithm}[definition]{Algorithm}
\theoremstyle{mlintuition}
\newtheorem*{intuition}{Intuition}
\newtheorem*{picturethis}{Picture this}
\theoremstyle{mlml}
\newtheorem*{mlnote}{Why this matters for ML}
\theoremstyle{mlpitfall}
\newtheorem*{pitfall}{Common pitfall}
\theoremstyle{mlremark}
\newtheorem*{remark}{Remark}
\theoremstyle{mlmotivation}
\newtheorem*{motivation}{Motivation}
\theoremstyle{mlnotation}
\newtheorem*{notation}{Notation}
\theoremstyle{mlexercise}
\newtheorem{exercisebase}{Exercise}[chapter]

\crefname{definition}{Definition}{Definitions}
\crefname{theorem}{Theorem}{Theorems}
\crefname{proposition}{Proposition}{Propositions}
\crefname{corollary}{Corollary}{Corollaries}
\crefname{example}{Example}{Examples}
\crefname{algorithm}{Algorithm}{Algorithms}
\crefname{exercisebase}{Exercise}{Exercises}
\crefname{figure}{Figure}{Figures}
\crefname{table}{Table}{Tables}
\crefname{equation}{Equation}{Equations}
\crefname{chapter}{Chapter}{Chapters}
\crefname{section}{Section}{Sections}
\crefname{subsection}{Section}{Sections}
\crefname{appendix}{Appendix}{Appendices}

\tcolorboxenvironment{definition}{mlbox=MLdef}
\tcolorboxenvironment{theorem}{mlbox=MLthm}
\tcolorboxenvironment{proposition}{mlbox=MLthm}
\tcolorboxenvironment{corollary}{mlbox=MLthm}
\tcolorboxenvironment{example}{mlbox=MLexample,colback=MLexample!3!white}
\tcolorboxenvironment{algorithm}{mlbox=MLalgorithm}
\tcolorboxenvironment{intuition}{mlbox=MLintuition,colback=MLintuition!7!white}
\tcolorboxenvironment{picturethis}{mlbox=MLintuition,colback=MLintuition!7!white}
\tcolorboxenvironment{mlnote}{mlbox=MLml}
\tcolorboxenvironment{pitfall}{mlbox=MLpitfall}
\tcolorboxenvironment{remark}{mlbox=MLremark,colback=MLremark!4!white}
\tcolorboxenvironment{motivation}{mlbox=MLmotivation}
\tcolorboxenvironment{notation}{mlbox=MLnotation}

% Proofs: amsthm's proof environment with a coloured QED box.
\renewcommand{\qedsymbol}{\textcolor{MLthm}{$\blacksquare$}}

% "Solution." label used inside worked examples.
\newcommand{\sol}{\par\smallskip\noindent\textit{Solution.}\enspace}

% A box for the learning goals at the start of each chapter.
\newtcolorbox{objectives}{
  enhanced,breakable,sharp corners,
  colback=MLchapter!4!white,colframe=MLchapter!60!white,boxrule=0.6pt,
  title={In this chapter you will learn to\dots},
  fonttitle=\bfseries\color{white},colbacktitle=MLchapter!85!white,
  left=8pt,right=8pt,top=5pt,bottom=5pt,before skip=10pt,after skip=12pt,
  parbox=false,
}

% A box for the summary at the end of each chapter.
\newtcolorbox{keytakeaways}[1][Key takeaways]{
  enhanced,breakable,sharp corners,
  colback=MLsummary!4!white,colframe=MLsummary!60!white,boxrule=0.6pt,
  title={#1},
  fonttitle=\bfseries\color{white},colbacktitle=MLsummary!85!white,
  left=8pt,right=8pt,top=5pt,bottom=5pt,before skip=10pt,after skip=10pt,
  parbox=false,
}

% A frame around a key formula, for use inside display maths:
%   \[ \keyeq{a^2+b^2=c^2} \]
% (The options are passed explicitly because tcolorbox resets the global
% "highlight math style" at the start of every box.)
\tcbset{mlkeyeq/.style={enhanced,sharp corners,colback=white,
  colframe=MLchapter!70!white,boxrule=0.8pt,left=5pt,right=5pt,
  top=3pt,bottom=3pt}}
\newcommand{\keyeq}[1]{\tcbhighmath[mlkeyeq]{#1}}

%----------------------------------------------------------------------
% New terms: bold, coloured, and automatically added to the index.
%   \term{gradient}               -> index entry "gradient"
%   \term[mean!arithmetic]{mean}  -> custom index entry
%----------------------------------------------------------------------
\newcommand{\term}[2][]{%
  \textbf{\color{MLterm}#2}%
  \ifblank{#1}{\index{#2}}{\index{#1}}}

%----------------------------------------------------------------------
% Exercises and automatically collected solutions
%
% Usage inside a chapter:
%
%   \sectionexercises
%   \begin{exercise}
%     Compute $3+4\times 2$.
%   \end{exercise}
%   \begin{solution}
%     Multiplication first: $3+8=11$.
%   \end{solution}
%
% The solution is not printed where it is written. Instead it is stored
% in the file \jobname.sol, together with the number of the exercise it
% belongs to, and \printsolutions (used in the solutions appendix)
% typesets all of them, grouped by chapter, with links in both
% directions. Solutions must not contain verbatim material or the
% character #.
%----------------------------------------------------------------------
\makeatletter
\newcounter{MLsolutioncount}
\newwrite\ML@solfile
\newif\ifML@solfileopen
\gdef\ML@lastsolchapter{}

\AtBeginDocument{%
  \immediate\openout\ML@solfile=\jobname.sol\relax
  \global\ML@solfileopentrue}

\newenvironment{exercise}[1][]{%
  \ifblank{#1}{\begin{exercisebase}}{\begin{exercisebase}[#1]}%
  \label{exercise:\theexercisebase}}%
  {\end{exercisebase}}
% The link to the solution sits in the bottom-right corner, in a strip
% of extra padding so that it never overlaps the text of the exercise.
\tcolorboxenvironment{exercise}{%
  mlbox=MLexercise,colback=MLexercise!3!white,bottom=12pt,
  overlay unbroken and last={%
    \node[anchor=south east,inner sep=3pt,font=\scriptsize]
      at (frame.south east)
      {\hyperref[solution:\theexercisebase]{\textcolor{MLexercise!75!white}%
        {Solution on p.\,\pageref*{solution:\theexercisebase}\,%
          $\rightarrow$}}};}}

\NewDocumentEnvironment{solution}{+b}{%
  \ifML@solfileopen
    \edef\ML@thischapter{\thechapter}%
    \ifx\ML@thischapter\ML@lastsolchapter\else
      \immediate\write\ML@solfile{%
        \string\MLSolutionsChapter{\thechapter}%
        {\unexpanded\expandafter{\MLcurrentchaptertitle}}}%
      \global\let\ML@lastsolchapter\ML@thischapter
    \fi
    \immediate\write\ML@solfile{%
      \string\MLSolution{\theexercisebase}{\unexpanded{#1}}}%
    \stepcounter{MLsolutioncount}%
  \fi}{}

% Heading printed before the exercises of a section.
\newcommand{\sectionexercises}{%
  \par\needspace{6\baselineskip}\medskip
  \noindent{\large\bfseries\color{MLexercise}Exercises for
    Section~\thesection}\par\nopagebreak\smallskip}

% Typesetting of the stored solutions.
\newcommand{\MLSolutionsChapter}[2]{%
  \section*{Chapter #1: #2}%
  \phantomsection
  \addcontentsline{toc}{section}{Chapter #1: #2}}
\newcommand{\MLSolution}[2]{%
  \par\needspace{3\baselineskip}\medskip\noindent
  \phantomsection\def\@currentlabel{#1}\label{solution:#1}%
  {\bfseries\hyperref[exercise:#1]{\textcolor{MLexercise}{Exercise~#1.}}}%
  \enspace #2\par}
\newcommand{\printsolutions}{%
  \ifML@solfileopen
    \immediate\closeout\ML@solfile
    \global\ML@solfileopenfalse
  \fi
  \InputIfFileExists{\jobname.sol}{}{%
    \par\textit{No solutions were recorded in this run.}\par}}
\makeatother

%----------------------------------------------------------------------
% Notation macros (see Appendix A for the full list)
%----------------------------------------------------------------------
% Number sets
\newcommand{\N}{\mathbb{N}}
\newcommand{\Z}{\mathbb{Z}}
\newcommand{\Q}{\mathbb{Q}}
\newcommand{\R}{\mathbb{R}}

% Vectors (bold lower case) and matrices (bold upper case)
\newcommand{\vect}[1]{\bm{#1}}
\newcommand{\mat}[1]{\bm{#1}}
\newcommand{\va}{\vect{a}}
\newcommand{\vb}{\vect{b}}
\newcommand{\vc}{\vect{c}}
\newcommand{\vd}{\vect{d}}
\newcommand{\ve}{\vect{e}}
\newcommand{\vg}{\vect{g}}
\newcommand{\vh}{\vect{h}}
\newcommand{\vk}{\vect{k}}
\newcommand{\vm}{\vect{m}}
\newcommand{\vp}{\vect{p}}
\newcommand{\vq}{\vect{q}}
\newcommand{\vr}{\vect{r}}
\newcommand{\vs}{\vect{s}}
\newcommand{\vu}{\vect{u}}
\newcommand{\vv}{\vect{v}}
\newcommand{\vw}{\vect{w}}
\newcommand{\vx}{\vect{x}}
\newcommand{\vy}{\vect{y}}
\newcommand{\vz}{\vect{z}}
\newcommand{\vzero}{\vect{0}}
\newcommand{\vone}{\vect{1}}
\newcommand{\vtheta}{\bm{\theta}}
\newcommand{\vmu}{\bm{\mu}}
\newcommand{\vdelta}{\bm{\delta}}
\newcommand{\vepsilon}{\bm{\varepsilon}}
\newcommand{\mA}{\mat{A}}
\newcommand{\mB}{\mat{B}}
\newcommand{\mC}{\mat{C}}
\newcommand{\mD}{\mat{D}}
\newcommand{\mH}{\mat{H}}
\newcommand{\mI}{\mat{I}}
\newcommand{\mM}{\mat{M}}
\newcommand{\mP}{\mat{P}}
\newcommand{\mQ}{\mat{Q}}
\newcommand{\mR}{\mat{R}}
\newcommand{\mS}{\mat{S}}
\newcommand{\mU}{\mat{U}}
\newcommand{\mV}{\mat{V}}
\newcommand{\mW}{\mat{W}}
\newcommand{\mX}{\mat{X}}
\newcommand{\mY}{\mat{Y}}
\newcommand{\mZero}{\mat{0}}
\newcommand{\mSigma}{\bm{\Sigma}}
\newcommand{\mLambda}{\bm{\Lambda}}

% Delimiters
\DeclarePairedDelimiter{\abs}{\lvert}{\rvert}
\DeclarePairedDelimiter{\norm}{\lVert}{\rVert}
\DeclarePairedDelimiter{\ceil}{\lceil}{\rceil}
\DeclarePairedDelimiter{\floor}{\lfloor}{\rfloor}

% Operators
\DeclareMathOperator*{\argmin}{arg\,min}
\DeclareMathOperator*{\argmax}{arg\,max}
\DeclareMathOperator{\Var}{Var}
\DeclareMathOperator{\Cov}{Cov}
\DeclareMathOperator{\Corr}{Corr}
\DeclareMathOperator{\Bias}{Bias}
\DeclareMathOperator{\MSE}{MSE}
\DeclareMathOperator{\tr}{tr}
\DeclareMathOperator{\rank}{rank}
\DeclareMathOperator{\diag}{diag}
\DeclareMathOperator{\sign}{sign}
\DeclareMathOperator{\spn}{span}
\DeclareMathOperator{\proj}{proj}
\DeclareMathOperator{\softmax}{softmax}
\DeclareMathOperator{\ReLU}{ReLU}
\DeclareMathOperator{\KL}{KL}
\newcommand{\E}{\mathbb{E}}
\newcommand{\dd}{\mathrm{d}}
\newcommand{\deriv}[2]{\frac{\dd #1}{\dd #2}}
\newcommand{\pderiv}[2]{\frac{\partial #1}{\partial #2}}
\newcommand{\iid}{\overset{\text{i.i.d.}}{\sim}}
